\section{Bases de Diseño PCI}

Las bases de diseño del sistema de protección contra incendios se fundamentan en la ocupación del proyecto, la clasificación del riesgo y las normas NFPA y NEC vigentes, registradas en el archivo \texttt{config/datos\_proyecto.tex}.

\subsection{Identificación del proyecto}

\begin{table}[htbp]
\centering
\caption{Identificación PCI}
\label{tab:pci_id}
\begin{tabularx}{\textwidth}{@{}lX@{}}
\toprule
\textbf{Parámetro} & \textbf{Valor} \\
\midrule
Ocupación / proceso & \occupancy \\
Clase de riesgo NFPA 13 & \hazardclass \\
Autoridad Competente (AHJ) & \ahj \\
Áreas clasificadas (NEC Art.\ 500) & \classifiedareas \\
Normas aplicables & \nfpacodes \\
\bottomrule
\end{tabularx}
\end{table}

\subsection{Condiciones ambientales del sitio}

\begin{table}[htbp]
\centering
\caption{Condiciones ambientales del sitio de instalación}
\label{tab:environmental_conditions}
\setcounter{itemcount}{0}
\begin{tabularx}{\textwidth}{@{}NXclc@{}}
\toprule
\multicolumn{1}{c}{\textbf{Item}} & \textbf{Parámetro Ambiental} & \textbf{Valor} & \textbf{Unidad} & \textbf{Referencia} \\
\midrule
& Temperatura ambiente máxima & 38{,}0 & °C & IDEAM \\
& Temperatura ambiente mínima & 15{,}0 & °C & IDEAM \\
& Humedad relativa máxima & 85{,}0 & \% & IDEAM \\
& Presión atmosférica & 101{,}3 & kPa & NOAA \\
& Altitud s.n.m. & 960 & msnm & IGAC \\
& Velocidad máxima del viento & 120{,}0 & km/h & Estación local \\
& Zona sísmica & -- & Intermedia & NSR-10 \\
\bottomrule
\end{tabularx}
\end{table}

\subsection{Bases hidráulicas (NFPA 13 / 14 / 20)}

\begin{table}[htbp]
\centering
\caption{Parámetros hidráulicos de diseño}
\label{tab:hidraulica_bases}
\setcounter{itemcount}{0}
\begin{tabularx}{\textwidth}{@{}NXlc@{}}
\toprule
\multicolumn{1}{c}{\textbf{Item}} & \textbf{Parámetro} & \textbf{Valor} & \textbf{Referencia} \\
\midrule
& Densidad de diseño & \densidadDiseno & NFPA 13 §19.2.3 \\
& Área de operación más remota & \areaOperacion & NFPA 13 §19.2 \\
& Tiempo mínimo de operación & \tiempoOperacion & NFPA 13 Tabla 19.2.3 \\
& Margen de seguridad hidráulico & \margenSeguridadHidro & Buenas prácticas \\
& Coeficiente Hazen-Williams (acero) & 120 & NFPA 13 Tabla 23.4.4.7.1 \\
& Bomba contra incendio & Según NFPA 20 & NFPA 20 \\
\bottomrule
\end{tabularx}
\end{table}

\subsection{Detección y alarma (NFPA 72)}

\begin{itemize}
\item Tipos de detectores seleccionados por área (calor, humo, llama UV/IR) según clase de riesgo y tipo de combustible presente.
\item Batería de respaldo dimensionada para 24 h en standby más 5 min en alarma (NFPA 72 §10.6.10), con factor de envejecimiento 1{,}20.
\item Caída de tensión en NAC/SLC verificada al dispositivo más alejado (NFPA 72 §10.6.8.2).
\item Nivel sonoro mínimo 15 dB por encima del ruido ambiente promedio (NFPA 72 §18.4).
\end{itemize}

\subsection{Extinción con agentes limpios (NFPA 2001, si aplica)}

\begin{table}[htbp]
\centering
\caption{Bases agente limpio}
\label{tab:agente_bases}
\begin{tabularx}{\textwidth}{@{}lX@{}}
\toprule
\textbf{Parámetro} & \textbf{Valor} \\
\midrule
Agente & \agenteLimpio \\
Concentración de diseño & \concentracionDiseno \\
Factor de seguridad & 1{,}20 (Clase B) / 1{,}30 (Clase A) \\
Tiempo de descarga & ≤ 10 s (halocarbonados) / ≤ 60 s (inertes) \\
Hold time & ≥ 10 min a altura de protección \\
\bottomrule
\end{tabularx}
\end{table}
